\begin{corollary}[uniform top-degree base change for surface fibres]
\hypertarget{tag-03DV}{}\label{wbII:cor:surface-top-base-change}
\StatementRecord{wbII:cor:surface-top-base-change}{proved}{03DV}{\SeriesRef{wbII:eq:W:5.17-U}, \SeriesRef{wbII:prop:analytic-base-change}; \cite[Chapter~III, \S4, Theorem~4.1, p.~124]{BS76push}}

\ Let $f:X\to B$ be a proper flat holomorphic map to a Riemann
surface whose full scheme-theoretic fibres have pure complex dimension two,
and let $E$ be a holomorphic vector bundle on $X$.  Then
\[
 R^qf_*E=0\quad(q>2)
 \tag{W.15}\hypertarget{tag-03FU}{}\label{wbII:eq:W:5.17-V}
\]
and, for every $b\in B$, including a point with nonreduced fibre, the
canonical map
\[
 (R^2f_*E)_b\otimes_{\mathcal O_{B,b}}\mathbb C(b)
 \xrightarrow{\ \sim\ }H^2(X_b,E|_{X_b})
 \tag{W.16}\hypertarget{tag-03FS}{}\label{wbII:eq:W:5.17-T}
\]
is an isomorphism.  No constancy of $h^2$, local freeness of $R^2f_*E$, or
reducedness of $X_b$ is required.
\end{corollary}



\begin{proof}
\begin{enumerate}[label={$\langle1\rangle$\arabic*.},leftmargin=*]
\item 
\noindent\textbf{Step 1.} On a sufficiently small neighbourhood $U$ of $b$, the finite-complex theorem
used in analytic cohomology and base change represents $Rf_*E|_U$ by a
bounded complex $K^\bullet$ of finite free $\mathcal O_U$-modules
\cite[Chapter~III, \S4, Theorem~4.1, p.~124]{BS76push}.  Flatness of
$E$ over $B$ identifies its derived fibre with the ordinary restriction to
the full scheme-theoretic fibre.  Over the discrete valuation ring
$A=\mathcal O_{B,b}$, the universal-coefficient sequence is
\[
0\longrightarrow H^q(K^\bullet)_b\otimes_A\mathbb C(b)
 \longrightarrow H^q(K^\bullet\otimes_A\mathbb C(b))
 \longrightarrow
 \operatorname{Tor}_1^A(H^{q+1}(K^\bullet)_b,\mathbb C(b))
 \longrightarrow0.
\tag{W.17}\hypertarget{tag-03FT}{}\label{wbII:eq:W:5.17-U}
\]
If some $R^qf_*E$ with $q>2$ were nonzero, choose the largest such $q$ and a
point of its support.  Nakayama's lemma makes the left term of \eqref{wbII:eq:W:5.17-U}
nonzero there, while maximality kills the Tor term.  This would give nonzero
coherent cohomology in degree $q>2$ on a compact two-dimensional complex
space, a contradiction.  Thus \eqref{wbII:eq:W:5.17-V} holds.  For $q=2$, the only
obstruction in \eqref{wbII:eq:W:5.17-U} is
$\operatorname{Tor}_1^A((R^3f_*E)_b,\mathbb C(b))$, which is zero by
\eqref{wbII:eq:W:5.17-V}, proving \eqref{wbII:eq:W:5.17-T}.  This is the explicit obstruction calculation
behind \cref{wbII:prop:analytic-base-change}..

\end{enumerate}
\end{proof}